\JOUChapterWithEnglish{浮选动力学}{Flotation Kinetics}
\label{chap:kinetics}

\JOUSectionWithEnglish{柱体背压的影响}{Influence of Back Pressure}

柱体背压直接影响气泡尺寸分布、矿浆停留时间以及循环流区的稳定性。背压较低时，气泡上升速度较快而停留时间不足；背压过高时，又可能增加能耗并削弱分选层内的相对稳定流场。因此，合理选择背压区间是保证动力学过程稳定推进的基础。

\JOUSectionWithEnglish{循环矿浆压力的影响}{Influence of Pressure of Circulating Pulp}

循环矿浆压力决定了旋流场强度和颗粒径向分离能力。随着压力上升，颗粒所受离心力增强，但若压力过大则会引起局部紊动增加和短路流现象。动力学分析通常需要在颗粒运动路径、界面吸附时间与能量消耗之间寻求折中关系。

\JOUSectionWithEnglish{循环矿浆量的影响}{Influence of Quantity of Circulating Pulp}

循环矿浆量既影响设备负荷，也改变柱内轴向速度分布与气泡携带能力。当循环矿浆量处于适中水平时，矿粒与气泡的有效接触频率更高，动力学速率常数可维持在较优区间；当负荷继续提高时，夹带效应与返混现象会削弱体系的选择性。
